\documentclass[12pt,a4paper]{article}

\usepackage[utf8]{inputenc}
\usepackage[greek,english]{babel}
\usepackage{alphabeta} 

\usepackage[pdftex]{graphicx}
\usepackage[top=1in, bottom=1in, left=1in, right=1in]{geometry}

\linespread{1.06}
\setlength{\parskip}{8pt plus2pt minus2pt}

\widowpenalty 10000
\clubpenalty 10000

\newcommand{\eat}[1]{}
\newcommand{\HRule}{\rule{\linewidth}{0.5mm}}

\usepackage[official]{eurosym}
\usepackage{enumitem}
\setlist{nolistsep,noitemsep}
\usepackage[hidelinks]{hyperref}
\usepackage{cite}
\usepackage{lipsum}


\begin{document}
\renewcommand{\abstractname}{Περίληψη}
\renewcommand{\contentsname}{Περιεχόμενα}
\renewcommand{\refname}{Αναφορές}

%===========================================================
\begin{titlepage}
\begin{center}

% Top 
\includegraphics[width=0.55\textwidth]{cut-logo-el}~\\[2cm]


% Title
\HRule \\[0.4cm]
{ \LARGE 
  \textbf{Αναφορά μαθήματος ΗΜ442}\\[0.4cm]
  \emph{Τίτλος αναφοράς}\\[0.4cm]
}
\HRule \\[1.5cm]



% Author
{ \large
  Όνομα Επίθετο \\[0.1cm]
  Μεταπτυχιακός Φοιτητής\\[0.1cm]
  \texttt{user@cut.ac.cy}
}

\vfill

%\textsc{\Large Cyprus University of Technology}\\[0.4cm]
\textsc{\large Τμήμα Ηλεκτρολόγων Μηχανικών, Μηχανικών\\Ηλεκτρονικών Υπολογιστών \& Πληροφορικής}\\[0.4cm]


% Bottom
{\large \selectlanguage{greek}\today}
 
\end{center}
\end{titlepage}


\begin{abstract}
Τα θεία βέλη στον στρατόν πετούσαν εννιά μέρες,
και την δεκάτην τον λαόν συγκάλεσε ο Πηλείδης,
ως η θεά τον διδάξεν η Ήρα η λευκοχέρα,
που εθλίβονταν τους Δαναούς να βλέπει πως πεθαίναν.
Και αφού συνάχθηκε ο λαός εις ένα μέρος όλος,
ο γοργοπόδης Αχιλλεύς σηκώθη και τους είπε·
«Ατρείδη, να γυρίσουμε, θαρρώ, θ' αναγκασθούμε
στα γονικά μας άπρακτοι, αν δεν πεθάνουμ' όλοι,
αφού μας φθέρνει λοιμική και πόλεμος αντάμα.
Λοιπόν ας ερωτήσομεν ή μάντιν ή ιερέα
ή ονειροκρίτην—έρχεται και τ' όνειρο απ' τον Δία —
να ειπεί γιατί εχόλωσε τόσο σ' εμάς ο Φοίβος·
μη κάποιο τάμα του έλειψε, μη του 'λειψ' εκατόμβη·
ίσως, αν του καούν αρνιά και ερίφια διαλεμένα,
θελήσει το θανατικό να διώξει από κοντά μας».
Αυτά είπε κι εκάθισε· και τότε ο Θεστορίδης
ο Κάλχας εσηκώθηκεν, ορνεοσκόπος πρώτος
που εγνώριζ' όλα μέλλοντα, παρόντα, περασμένα
και οδήγησε στην Ίλιον των Αχαιών τα πλοία
μ' αυτό το πνεύμα μαντικό που του 'χε δώσει ο Φοίβος.
\addtocontents{toc}{\protect\thispagestyle{empty}}
\end{abstract}

\newpage



%===========================================================
\tableofcontents
\addtocontents{toc}{\protect\thispagestyle{empty}}
\newpage
\setcounter{page}{1}

%===========================================================
%===========================================================
\section{Εισαγωγή}\label{sec:intro}

Σ' αυτούς καλοπροαίρετα τότε ομιλούσ' εκείνος:
«Με προσκαλείς, διίφιλε Πηλείδη, να εξηγήσω,
πώς εγεννήθηκε ο θυμός του μακροβόλου Φοίβου·
θέλει το ειπώ· μόνον εσύ στοχάσου και όμοσέ μου
να με βοηθήσεις πρόθυμα με λόγον και με χέρι,
ότι θ' ανάψω την χολήν ανδρός που των Αργειών
δεσπόζει και ολ' οι Αχαιοί του είναι υποταγμένοι·
όταν θυμώσει στον μικρόν, νικά ο βασιλέας·
ότι αν χωνεύσει την χολήν σ' εκείνην την ημέραν,
όμως το πάθος άσπονδο στα στήθη μέσα τρέφει
να ξεθυμάνει στο εξής· και σκέψου αν θα με σώσεις». \cite{iliada}

\section{Βιβλιογραφική αναφορά}\label{sec:lit-rev}

Και ο γοργοπόδης προς αυτόν Πηλείδης αποκρίθη:
«Άφοβα λέγε τον χρησμόν όποιον ηξεύρει ο νους σου·
ότι, μα τον Απόλλωνα, που τες ευχές ακούει,
Κάλχα, και συ των Δαναών προσφέρεις τους χρησμούς του,
όσο εγώ ζω κι εδώ στην γην βλέπω το φως του ηλίου,
βαρύ κανείς επάνω σου το χέρι δεν θα βάλει
των Δαναών όλων κανείς, και μήτε ο Αγαμέμνων
που σήμερα των Αχαιών καυχάται ότ' είναι ο πρώτος».\cite{latexcompanion}

\section{Μεθοδολογία}\label{sec:meth}

Και ο μάντις ο ακατάκριτος επήρε θάρρος κι είπε:
«Τάμα ποσώς δεν του 'λειψε, μήτ' εκατόμβη, αλλ' είναι
ο ιερέας αφορμή, που αψήφησ' ο Ατρείδης·
την κόρη δεν απόλυσε, τα λύτρα δεν εδέχθη,
ιδού γιατί μας έθλιψε και θα μας θλίψει ο Φοίβος.
Ουδ' απ' τους Δαναούς ποτέ την λοιμική θα διώξει
πριν δοθεί οπίσω του πατρός η λαμπρομάτα κόρη
άλυτρη, ανεξαγόραστη και αγίαν εκατόμβην
στον Χρύσην αποστείλομεν· τότ' ίσως ίλεως γίνει».

\section{Αποτελέσματα}\label{sec:res}

Αυτά είπε κι εκάθισε· σηκώθη ευθύς ο ήρως
πολλών κυρίαρχος λαών, ο Ατρείδης Αγαμέμνων
φαρμακωμένος· και η χολή τα μαύρα σωθικά του
πλημμύριζ' όλα, και άστραφταν τα μάτια του ωσάν φλόγες.
Με βλέμμα κακοσήμαντο στον Κάλχαντα είπε πρώτα:
«Μάντι κακών, όχι, ποτέ πρόσχαρό τι δεν μου 'πες,
και ο νους σου πάντοτε αγαπά κακά να προμαντεύει·
λόγον δεν είπες συ ποτέ καλόν ούτ' έχεις πράξει.
Και τώρα εδώ στους Δαναούς χρησμολογείς και λέγεις,
οπώς για τούτο συμφορές τους δίδει ο μακροβόλος,
ότι την πλούσια ξαγορά της θυγατρός του Χρύση
δεν δέχθηκα· ναι, θέλω εγώ καλύτερα την κόρη
σπίτι μου, αφού την προτιμώ της νυμφευτής μου ακόμα
της Κλυταιμνήστρας και ποσώς κατώτερη δεν είναι
στην κλάση, στο ανάστημα, στη γνώμη και στα έργα.
Και όμως αν συμφέρει αυτό, θε να την αποδώσω·
το καλό θέλω του λαού, ποτέ τον όλεθρό του·
αλλά δώρο ετοιμάσετε σ' εμένα ευθύς, τι μόνος
εγώ δεν πρέπει αδώρητος να μείνω των Αργείων
και όλοι το βλέπετε ότι αλλού το δώρο μου πηγαίνει».


%===========================================================
%===========================================================
\addcontentsline{toc}{section}{Αναφορές}
\bibliographystyle{ieeetr}
\bibliography{refs}


\end{document} 
